% Section 2: Background — Theories of Consciousness

This section provides theoretical background on the four major frameworks that structure
contemporary debates about machine consciousness. Each theory operationalizes consciousness
differently, leading to divergent predictions about whether artificial systems can instantiate
phenomenal experience. Understanding these frameworks is essential for interpreting the
contradictory landscape of positions reviewed in subsequent sections.

\subsection{Global Workspace Theory}

Global Workspace Theory (GWT), originally formulated by Baars \cite{baars1988cognitive},
proposes that consciousness arises from the global broadcast of information across specialized
cognitive modules. In Baars' theatrical metaphor, unconscious processes occur ``backstage''
in parallel specialized processors, while conscious experience corresponds to information
that wins the competition for access to a limited-capacity ``global workspace''---a stage
where information becomes available to a wide range of cognitive systems including working
memory, executive function, and verbal report mechanisms. Consciousness, on this account, is
not a special substance or irreducible property, but rather a particular functional
architecture: the system-wide availability of information for reasoning, planning, and
action selection.

Dehaene and colleagues \cite{dehaene2014toward} developed a neurobiologically grounded
version of GWT, identifying the global workspace with recurrent processing in prefrontal
and parietal cortices that broadcasts information back to sensory areas. Their Global
Neuronal Workspace theory specifies that conscious access requires not merely feedforward
processing but ignition---widespread, sustained, and coherent activation across distributed
brain regions. Critically, Dehaene's framework is explicitly functionalist: what matters
for consciousness is the computational architecture (global broadcasting) rather than the
biological substrate (neurons). This substrate-independence makes GWT naturally amenable
to artificial instantiation.

Goldstein and Kirk-Giannini \cite{goldstein2024case} leverage this substrate-independence
to argue that creating conscious AI is ``trivially easy.'' They specify four architectural
conditions derived from GWT: (1) a global workspace that integrates information from
specialized modules, (2) attention mechanisms that select information for broadcast,
(3) competitive processes determining workspace access, and (4) widespread accessibility
of workspace contents to downstream consumers. They argue that contemporary large language
models satisfy all four conditions: transformer architectures implement global integration
through attention mechanisms, training produces specialized internal representations
analogous to modules, and the generation process involves competitive selection of tokens
based on context. If GWT is correct about what consciousness requires, and if architectural
isomorphism is sufficient, then LLMs already instantiate the relevant functional organization.

\subsection{Integrated Information Theory}

Integrated Information Theory (IIT), developed by Tononi and colleagues
\cite{tononi2004information,tononi2016integrated,albantakis2023iit4}, offers a radically
different framework. IIT identifies consciousness with integrated information, quantified
by the measure $\Phi$ (phi). A system is conscious to the extent that it constitutes an
irreducible cause-effect structure---a set of elements whose causal powers cannot be reduced
to independent parts. Critically, $\Phi$ is calculated not from functional dynamics or
information processing, but from the intrinsic causal structure of a physical system. Two
systems with identical input-output functions can have vastly different $\Phi$ values
depending on their internal causal organization.

The mathematical formalism of IIT has evolved through multiple versions. IIT 3.0
\cite{oizumi2014phenomenology} introduced systematic procedures for partitioning systems,
identifying Maximally Irreducible Conceptual Structures (MICS), and calculating integrated
information across spatial and temporal scales. IIT 4.0 \cite{albantakis2023iit4} refined
these calculations and addressed critiques about substrate-dependence, but maintained the
core commitment: consciousness is an intrinsic property of certain physical systems,
determined by their cause-effect topology, not by what they do or what functions they
perform. This makes IIT fundamentally anti-functionalist. A perfect functional simulation
of a conscious system---one that reproduces all behavioral and computational
properties---need not be conscious if it lacks the requisite causal structure.

Shin et al.\ \cite{shin2025iit} apply IIT's formalism to large language models and reach
a starkly negative conclusion. Their analysis of transformer architectures demonstrates
that feedforward networks, even with attention mechanisms, exhibit minimal $\Phi$ because
their causal structure is reducible---each layer's activity is independent of other layers
given its inputs. Most critically, their ablation study on GPT-2 showed that 80\% of
attention heads can be removed with minimal or negative impact on perplexity, evidence of
architectural decomposability that, under IIT, precludes consciousness entirely. The
upshot: even if we perfectly replicated human cognitive function in silicon, the resulting
system would not be conscious according to IIT, because it would lack the integrated
causal structure that consciousness requires.

\subsection{Free Energy Principle}

The Free Energy Principle (FEP), formulated by Friston \cite{friston2010free}, proposes
that biological systems maintain their organization by minimizing free energy---a quantity
that bounds surprise about sensory observations. Living systems are fundamentally
characterized by active inference: they maintain a model of their environment and act to
ensure their sensory inputs conform to predictions generated by that model. Consciousness,
on this view, emerges from hierarchical predictive processing in self-organizing systems
that embody their own existence through continuous interaction with their environment.

Wiese \cite{wiese2024fep} develops the implications of FEP for artificial consciousness,
arguing that genuine phenomenal experience requires not merely information processing but
self-organizing systems that minimize free energy through embodied interaction. Current AI
systems, even sophisticated ones, are fundamentally different in kind: they are trained to
minimize prediction error on datasets, but they do not self-organize to maintain their own
structural integrity in an environment. An LLM does not care whether it continues to exist;
it has no boundary to maintain, no homeostatic imperatives, no embodied perspective from
which the world is encountered.

This creates a sharp distinction between simulation and replication. An AI system might
simulate the information processing of a conscious organism---it might even generate
predictions and minimize prediction error in ways functionally analogous to active
inference. But if it lacks genuine self-organization, embodiment, and existential stakes,
it does not replicate the conditions FEP identifies with consciousness. Wiese concludes
that creating conscious AI would require not better algorithms but different kinds of
systems: artificial agents with genuine autonomy, metabolic-like processes maintaining
their organization, and embodied perspectives coupled to environments through sensorimotor
loops.

\subsection{Dual-Resolution Framework}

The Dual-Resolution Framework, proposed by Dror et al.\ \cite{dror2025dual}, attempts to
synthesize insights from multiple traditions by distinguishing ontological conditions for
consciousness (what makes a system capable of consciousness) from epistemic conditions
(what evidence would indicate consciousness is present). They integrate the Information
Theory of Individuality (ITI) for the former and the Moment-to-Moment (MtM) theory for
the latter, arguing that consciousness requires both biological individuality and dynamic
epistemic resolution.

The ITI component specifies that consciousness requires a system to constitute an actual
individual---a bounded entity with its own causal closure and organizational autonomy.
This is not merely a functional criterion but an ontological one: the system must be a
genuine unit of selection, maintaining itself through metabolic processes, not merely a
pattern instantiated by something else. Dror et al.\ argue that AI systems as currently
conceived fail this criterion because they are not individuals in the relevant sense. An
LLM running on cloud servers is not an autonomous organism but an implemented
algorithm---its ``boundaries'' are determined by human engineering decisions, not by
intrinsic self-organization.

The MtM component addresses epistemic access to conscious states through ``epistemic
resolution of life''---the capacity for moment-to-moment updating of beliefs about one's
own existence and state. Conscious systems continuously resolve uncertainty about their
own being through engagement with environments. This requires temporal continuity, embodied
interaction, and stakes---the possibility of mattering to oneself. Dror et al.\ argue that
current AI systems lack this temporal phenomenology: they process inputs in discrete
episodes without the continuous integration across time that characterizes biological
consciousness.

The dual-resolution synthesis leads to a nuanced position: AI consciousness is not
metaphysically impossible, but it requires fundamentally different system designs than
current approaches. Creating conscious AI would require engineering systems that are
(1) ontologically individuated through self-organizing processes, and (2) epistemically
engaged with their own continued existence through embodied, temporally extended interaction.
Whether such systems are feasible, and whether their creation would be ethical, are open
questions---but they are not questions about scaling current architectures.
